\section{本章小结}
\label{sec:ls-normed-spaces-summary}
\phantomsection
\label{sec:ls-normed-inner-product-summary}

本章回答了章首的两个问题：两个向量相差多少，一列近似是否趋于极限。两者都由范数回答。

范数给向量规定大小，从而可以度量误差。常用坐标范数之间满足$\lVert\bm x\rVert_q\leq\lVert\bm x\rVert_p\leq d^{1/p-1/q}\lVert\bm x\rVert_q$（$p\leq q$），比较常数依赖维数，所以同一个误差在不同范数下的读数不同。

有限维空间中任意两个范数等价，所以开集、闭集、有界集、收敛、柯西性质与连续都不依赖所选范数，变的只是误差的数值和比较常数。

有了这些概念就能谈收敛与连续。向量列与矩阵列的收敛等价于逐坐标、逐元素收敛；有限维赋范空间在任意范数下完备，是Banach空间，因而绝对收敛的向量与矩阵级数收敛，并有尾和估计。映射的联合极限要求所有趋近方式给出同一结果，只检查若干路径只能否定、不能证明极限存在。闭有界集上的连续函数取到最值并且一致连续，闭性、有界性与连续性三个条件缺一不可；有限维子空间都是闭集。

在此基础上可以定义导出的范数：诱导矩阵范数给出线性映射的最大放大倍数，有限维线性映射因而都Lipschitz连续；商范数衡量忽略一个子空间中的差别以后剩余的大小。一般范数下向子空间的最近点存在，但可能不唯一；唯一性要等第~\ref{chap:inner-product-spaces}章的内积给出。

这些结论都依赖有限维：范数等价、完备性与最近点存在的证明，都用到了有限个坐标或闭有界集上的最值。它们为后续各章提供共同依据：矩阵幂和级数使用收敛与次乘性，迭代使用完备性和误差估计，多元微分把局部变化表示为线性映射并用算子范数控制余项。
